\section{Local dark matter distribution}
\label{sec: distributions}

Computations of DM direct detection event rates strongly depend on the assumptions made for the DM  distribution in the Solar neighborhood. In sections \ref{sec: DM density} and \ref{sec: velocity distributions}, we present the DM density and velocity distribution in the Solar neighborhood extracted from the simulated MW-LMC analogues, and discuss the effect of the LMC on the results.


For each possible Sun's position (and velocity) which matches the observed Sun-LMC geometry, we consequently have the orientation of the $(x_g, y_g, z_g)$ axes of the Galactic reference frame defined in section~\ref{sec: Sun-LMC Geometry}. We then transform the positions and velocities of the simulation particles to  this Galactic reference frame.  

\begin{figure}[t]
    \centering
    \includegraphics[scale=0.45]{graphics/roi_new_smallsun-cropped.pdf}
    \caption{Diagram of the Solar region (blue) chosen as the overlap of the volume of a cone projected from the Galactic center with an opening angle of $\pi/4$ radians and its axis fixed on the position of the Sun, with the volume enclosed between two concentric spheres with radii of 6 and 10~kpc from the Galactic center. For illustration, the Sun is placed on the cone's axis at a Galactocentric distance of 8 kpc.
    }
    \label{fig: solar region}
\end{figure}

To define the \emph{Solar region}, with the Sun at a galactocentric distance of $\sim 8$~kpc, we first consider the region enclosed within a spherical shell between 6 to 10 kpc from the Galactic center of the MW analogue. We then consider a cone with an opening angle of $\pi/4$ radians, its vertex at the Galactic center, and its axis aligned with the position of the Sun as obtained from the procedure discussed in section~\ref{sec: Sun-LMC Geometry}. The overlap of the spherical shell and the cone constitutes the Solar region, shown as the shaded blue region in figure~\ref{fig: solar region}. The size of the Solar region is chosen to be large enough to include several thousand DM particles, 
and small enough to retain sensitivity to the best fit Sun's position. In sections~\ref{sec: DM density} and \ref{sec: Variation due to Sun-LMC Geometry}, we discuss the impact of changing the size of the Solar region on the local DM density, the percentage of the DM particles originating from the LMC in the Solar region, and the high speed tails of the halo integrals.


Since the set of allowed and best fit Sun's positions we find using the procedure described in section~\ref{sec: Sun-LMC Geometry} vary for each halo and snapshot, the Solar region is  different for each MW analogue and each snapshot.

\begin{table}[t!]
    \centering
    \begin{tabular}{|c|c|c|c|c|c|c|}
      \hline
      Halo ID  & $N_{\rm MW}$ & $N_{\rm LMC}$ & $\rho_{\chi}$~[GeV$/$cm$^3$] & $\kappa_{\rm LMC}$~[\%] & $\kappa_{\rm LMC}$ Range~[\%] & $v_{\rm esc}^{\rm det}$~[km/s] \\
       \hline
       1  & 7,760 & 11 & 0.21 & 0.14 & $[0.14-0.21]$ & 651  \\
       2  & 8,581 & 55 & 0.23 & 0.64 & $[0.53-0.65]$ & 720 \\
       3 & 11,621 & 3 & 0.35 & 0.026 & $[0.025-0.028]$ & 714 \\
       4  & 12,483 & 12 & 0.34 & 0.096 & $[0.088-0.12]$ & 737 \\
       5  & 8,669 & 136 & 0.24 & 1.5 & $[1.4-1.6]$ & 707 \\
       6  & 13,290 & 5 &  0.38 & 0.038 & $[0.029-0.046]$ &  734 \\
       7  & 18,467 & 6 &  0.53 & 0.032 & $[0.032-0.034]$ & 766 \\
       8  & 12,949 & 1 & 0.38 & 0.0077 & $[0.0077-0.0082]$ & 712 \\
       9  & 11,892 & 12 & 0.36 & 0.10 & $[0.069-0.13]$ & 715 \\
       10 & 12,405 & 361 & 0.39 & 2.8 & $[2.8-3.1]$ & 791\\
       11  & 14,132 & 4 & 0.43 & 0.028 & $[0.021-0.039]$ & 758 \\
       12  & 16,427 & 28 & 0.53 & 0.17 & $[0.17-0.21]$ & 783 \\
       13 & 10,814 & 254 &  0.34 & 2.3 & $[2.3-3.0]$ & 831 \\
       14  & 20,001 & 52 & 0.60 & 0.26 & $[0.26-0.31]$ & 776 \\
       15 & 10,641 & 128 &  0.32 & 1.2 & $[0.81-1.3]$ & 819 \\
      \hline
    \end{tabular}%}
\caption{Various quantities for the 15 MW-LMC analogues in the Solar region, at the simulation snapshot closest to the LMC's pericenter approach: halo ID, the number of native DM particles of the MW, $N_{\rm MW}$, the number of DM particles originating from the LMC, $N_{\rm LMC}$, the local DM density, $\rho_\chi$, the  percentage of the DM particles originating from the LMC in the Solar region, $\kappa_{\rm LMC}$, for the best fit Sun's position, the range that $\kappa_{\rm LMC}$ can span across the different allowed Sun's positions, and the local escape speed from the MW in the detector rest frame, $v_{\rm esc}^{\rm det}$. All columns, except the 6th, list the quantities for the best fit Sun's position.} 
    \label{tab:Localrho}
  \end{table}  

The number of the native DM particles of the MW, $N_{\rm MW}$, and the number of the DM particles originating from the LMC, $N_{\rm LMC}$, in the Solar region for the best fit Sun's position are listed in table~\ref{tab:Localrho} for the 15 MW-LMC analogues at the snapshot closest to LMC's first pericenter approach. While there are $[7,760-20,001]$  DM particles from the MW in the Solar region, the number of DM particles originating from the LMC in the Solar region is in the range of $[1 - 361]$. Due to this limited number of LMC particles in the Solar region, we are not sensitive
to the variation of the distribution of DM particles from the LMC within our defined Solar region. The low number of DM particles originating from the LMC is, therefore, a limitation of the current cosmological simulations as compared to idealized simulations, which can achieve a better resolution. Nevertheless, due to their high relative velocities with respect to the Sun, the DM particles from the LMC are more numerous compared to the high speed DM particles of the MW, and can significantly affect the high speed tails of the local DM velocity distribution  (as discussed below in section~\ref{sec: velocity distributions}). Therefore, the low value of $N_{\rm LMC}$ is not a major concern for the validity of our results. 





\subsection{Local dark matter density}
\label{sec: DM density}


We first extract the local DM density, $\rho_\chi$, in the Solar region for the best fit Sun's position for the 15 MW-LMC analogue systems in Auriga at the snapshot closest to LMC's first pericenter approach. The results are given in table~\ref{tab:Localrho}. 
The local DM density is in the range of $\rho_\chi=[0.21-0.60]$~GeV/cm$^3$. This agrees  with the values obtained previously for the local DM density of MW-like halos in the EAGLE and APOSTLE~\cite{Bozorgnia:2016ogo}, and Auriga~\cite{Bozorgnia:2019mjk} simulations. It also agrees well with the local~\cite{Salucci:2010qr, Smith:2011fs, Bovy:2012tw, Garbari:2012ff, Zhang:2012rsb, Bovy:2013raa, 2018A&A...615A..99H, Buch:2018qdr} and global~\cite{McMillan:2011wd, Catena:2009mf, Weber:2009pt, Iocco:2011jz, Nesti:2013uwa, Sofue:2015xpa, Pato:2015dua, deSalas:2019pee} estimates from observations. The large range of local DM densities obtained from simulations is due to halo-to-halo variations and depends on halo properties such as mass (in our case the simulated halos have a mass to within less than a factor of 2 of that estimated for the MW halo~\cite{Callingham:2018vcf}), concentration, formation history, and mass of the stellar disk.

Next, we extract the percentage of the DM particles in the Solar region originating from the LMC analogue, $\kappa_{\rm LMC}$, at the snapshot closest to LMC's first pericenter approach. We consider a DM particle to have originated from the LMC analogue if it is bound to the LMC at infall as identified by the SUBFIND algorithm, and its distance from the center of the LMC at infall is less than twice the virial radius of the LMC at infall\footnote{The  majority of the  DM particles from the LMC in the Solar neighbourhood are typically found 
well within the virial radius of the LMC at infall and are therefore highly bound to the LMC at infall. Thus, our results are robust with respect to  the way we select the DM particles that have originated from the LMC analogue.}. $\kappa_{\rm LMC}$ is defined as the ratio of the number of DM particles originating from the LMC analogue in the Solar region and the total number of DM particles in the Solar region, multiplied by 100 to obtain the percentage. For the 15 MW-LMC analogues, $\kappa_{\rm LMC}$ in the Solar region for the best fit Sun's position is in the range of $[0.0077-2.8]$\%, as listed in table~\ref{tab:Localrho}. 
In the fourth column of the table,  we present the range that $\kappa_{\rm LMC}$ varies for each halo due to the different allowed Sun's positions. 





To investigate the reason for the halo-to-halo variation in $\kappa_{\rm LMC}$ and $N_{\rm LMC}$, in figure~\ref{fig: Correlations} we present the variation of these parameters with $M_{\rm Infall}^{\rm LMC}/M_{200}^{\rm MW}$ and $M_{\rm Infall}^{\rm LMC}$, respectively. The point sizes increase with the distance of the LMC analogues from host at pericenter. The left panel of the figure shows that in general, systems with a larger LMC to MW halo mass ratio also have a larger percentage of LMC particles in the Solar region in most cases. However, the two parameters are not tightly correlated. In particular, systems with similar $M_{\rm Infall}^{\rm LMC}/M_{200}^{\rm MW}$ can still show a large variation in $\kappa_{\rm LMC}$. This is mainly due to the variation in the distance of the LMC analogues from host at pericenter, $r_{\rm LMC}$, for these systems. A larger $r_{\rm LMC}$ translates to smaller $\kappa_{\rm LMC}$ for systems with similar LMC to MW mass ratio. Similarly, the right panel of the figure shows a degree of correlation between $N_{\rm LMC}$ and $M_{\rm Infall}^{\rm LMC}$, while there exists a degree of inverse correlation between $N_{\rm LMC}$ and $r_{\rm LMC}$ for systems with similar $M_{\rm Infall}^{\rm LMC}$. 



\begin{figure}[t]
    \centering
    \includegraphics[width=\textwidth]{graphics/fig_correlations.pdf}
    \caption{
    The correlation between $\kappa_{\rm LMC}$ and $M_{\rm Infall}^{\rm LMC}/M_{200}^{\rm MW}$ (left), and $N_{\rm LMC}$ and $M_{\rm Infall}^{\rm LMC}$ (right) for the 15 MW-LMC analogues. $\kappa_{\rm LMC}$ and $N_{\rm LMC}$ are given in the Solar region for the best fit Sun's position at the simulation snapshot closest to the LMC's pericenter approach. The sizes of points increase with the distance of the LMC analogues from host at pericenter.
    }
    \label{fig: Correlations}
\end{figure}


We have also checked how $\rho_\chi$ and $\kappa_{\rm LMC}$ vary if we change the size of our defined Solar region. In particular, for the re-simulated halo 13 at the present day snapshot, decreasing the opening angle of the cone from $\pi/4$ to $\pi/6$ while keeping the spherical shell width the same, cuts $N_{\rm LMC}$ and $N_{\rm MW}$ by half, decreases  $\rho_\chi$ by $\sim 30\%$,  and increases $\kappa_{\rm LMC}$ by $\sim 20\%$, compared to the original Solar region. Decreasing the shell width from $6-10$~kpc to $7-9$~kpc while keeping the opening angle of the cone the same has a similar effect on $N_{\rm LMC}$ and $N_{\rm MW}$, but leads to an increase of $\sim 2\%$ in $\rho_\chi$  and $\sim 10\%$ in $\kappa_{\rm LMC}$. Decreasing both the opening angle of the cone to $\pi/6$ and the shell width to $7-9$~kpc, reduces $N_{\rm LMC}$ to $1/3$ and $N_{\rm MW}$ to $1/4$ of their original values, decreases $\rho_\chi$ by $\sim 25\%$,  and increases  $\kappa_{\rm LMC}$ by $\sim 35\%$. These changes are smaller than the halo-to-halo variation in these parameters, as it can be seen from table~\ref{tab:Localrho}.
  

\subsection{Dark matter velocity distributions}
\label{sec: velocity distributions}



Next we extract the DM speed distributions in the Solar region in the Galactic reference frame. For each halo, the  velocity vectors of the DM particles are specified with respect to the halo center. The normalized DM  speed distribution, $f(v)$, is given by
\be
f(v) = v^2 \int d\Omega_{\mathbf{v}}\Tilde{f}(\mathbf{v}) \; ,
\ee
where  $d\Omega_{\mathbf{v}}$ is an infinitesimal solid angle around the direction $\mathbf{v}$, and  $\Tilde{f}(\mathbf{v})$ is the normalized DM velocity distribution such that $\int dv f(v)=\int d^3v \Tilde{f}(\mathbf{v}) = 1$.




In the SHM, the local circular speed of the MW is usually set to 220~{\rm km/s}. To compare the local DM speed distributions of different halos, we scale the DM speeds in the Solar region for each halo by $(220~{\rm km/s})/v_c$, where $v_c$ is the local circular speed 
computed from the total mass enclosed within a sphere of radius 8~kpc for each halo. Moreover, we choose an optimal speed bin size of $25 \, \rm km/s$ to compute the DM speed distributions from the simulations. This bin size ensures that there are enough particles in each speed bin such that the statistical noise in the data points remains small, without smearing out any possible features in the DM speed distributions.


In figure~\ref{fig: fv lmc candidates} we present the DM speed distributions in the Galactic rest frame for four MW-LMC analogues in the Solar region specified by their best fit Sun's position, for the  snapshot closest to the LMC's pericenter approach. The speed distribution of the total DM particles (native to the MW\footnote{Notice that the DM particles native to the MW in this simulation snapshot are under the influence of the LMC and are different from the DM particles belonging to an isolated MW.}  or originating from the LMC) in the Solar region is shown as black shaded bands (specifying the $1\sigma$ Poisson errors), while the distribution of the DM particles native to the MW is shown in red. The blue shaded bands show the speed distributions of the DM particles originating from the LMC in the Solar region, scaled down by a factor of 10 for better visualization. The speed distribution of the total DM particles and those native to the MW are both normalized to 1. The percentage of the DM particles in the Solar region originating from the LMC is also specified in the top left corner of each panel. The panels below the speed distribution plots
show the ratio of the speed distribution of the total DM particles and the MW-only distribution.

\begin{figure}[t]
    \centering
    \includegraphics[width=0.8\textwidth]{graphics/fv_panel_ratio_1.pdf}
    \caption{
    DM speed distributions in the Galactic rest frame in the Solar region for the best fit Sun's position  for four representative MW-LMC analogues: halo 2 (top left), halo 6 (top right), halo 13 (bottom left) and halo 15 (bottom right), for the  snapshot closest to the LMC's pericenter approach. The distributions of the DM particles originating from the MW+LMC, MW only, and LMC only are shown as black, red, and blue shaded bands specifying the $1\sigma$ Poisson errors, respectively. The LMC-only distribution has been scaled down by a factor of 10 for better visualization. The percentage of the DM particles originating from the LMC in the Solar region, $\kappa_{\rm LMC}$, is also specified on each panel in the upper left corner. The panels below the speed distribution plots show the ratio between the MW+LMC  and the MW-only distributions.
    }
    \label{fig: fv lmc candidates}
\end{figure}

Among the 15 MW-LMC analogues, the four halos presented in figure~\ref{fig: fv lmc candidates} are representative of the differences seen in the local speed distributions of the DM particles originating from the MW only, the LMC only and the combined MW+LMC. Halo 2 (top left) has an intermediate percentage of DM particles originating from the LMC in the Solar region ($\kappa_{\rm LMC}=0.64$\%). It also has a sharply peaked speed distribution, leading to noticeable differences between the tails of the MW+LMC and MW-only speed distributions, with their ratio reaching values greater than 2 in the tail. Halo 6  (top right) is an example of a halo for which even a small fraction of DM particles in the Solar region originating from the LMC  ($\kappa_{\rm LMC}=0.038$\%) can lead to   differences in the tail of its DM speed distribution, as seen from the ratio plot. Halo 13 (bottom left) has a high  fraction of DM particles originating from the LMC ($\kappa_{\rm LMC}=2.3$\%) with a broad speed distribution, leading to  mild differences  between the MW+LMC and MW-only speed distributions  across a large range of speeds. The ratio of the two distributions reaches similar values in halo 6 and halo 13, despite halo 13 having a $\kappa_{\rm LMC}$ which is $\sim 60$ times larger than halo 6. Finally, halo 15 (bottom right) with $\kappa_{\rm LMC}=1.2$\%, 
shows a large variation  between the MW+LMC and MW-only speed distributions in the high speed tail, with their ratio approaching 4. 



In general, the speed distribution of DM particles originating from the LMC is found to peak at the high speed tail ($\gtrsim 500$~km/s with respect to the Galactic center) of the speed distribution of DM particles originating from the MW. This leads to variations in the tail of the MW+LMC speed distribution as compared to the MW-only distribution, although the degree to which the distributions vary is subject to large halo-to-halo scatter. The particular shape and width of the LMC's speed distribution in the Solar region for each MW analogue can  affect the variations in the tail of the MW+LMC distribution. For example, halos with an even larger $\kappa_{\rm LMC}$ (as listed in table~\ref{tab:Localrho}), do not necessarily show significant differences in their $f(v)$ with and without the LMC particles. 






In order to explore further the impact of the LMC on the local DM distribution during its orbit, we next focus on halo 13, where we rerun the simulations with finer snapshots close to the LMC's pericenter approach, as discussed in section~\ref{sec: selection criteria}. 
In figure \ref{fig: fv panel plot halo 49} we present the local DM speed distributions in the Galactic rest frame for halo 13 for the four snapshots representing different times in the LMC's orbit of the MW analogue (given in table~\ref{tab:snapshots}). The local speed distributions of the DM particles originating from the MW only (red), the LMC only (blue), and the MW+LMC (black) are shown. The distributions are presented in the Solar region for the best fit Sun's position for all snapshots 
except for the isolated MW, for which there is no LMC analogue and the best fit Sun's position cannot be defined. Hence, for the isolated MW the DM distribution is extracted in a spherical shell with radii between 6 to 10~kpc from the Galactic center. $\kappa_{\rm LMC}$ is also specified in each panel. The panels below the speed distribution plots show the ratio of the MW+LMC and the MW-only distributions, for all snapshots other than the isolated MW snapshot. 


\begin{figure}[t]
    \centering
    \includegraphics[width=0.8\textwidth]{graphics/fv_panel_ratio_2.pdf}
    \caption{
    Local DM speed distribution in the Galactic rest frame for halo 13 for four representative snapshots: isolated MW (top left), LMC's pericenter approach (top right), present day MW-LMC (bottom left), and future MW-LMC (bottom right).
    The speed distributions of the DM particles originating from the MW+LMC, MW only, and LMC only are shown in black, red, and blue shaded bands representing the $1\sigma$ Poisson errors, respectively. The LMC-only distribution has been scaled down by a factor of 10. The distributions are presented in the Solar region for the best fit Sun's position, except for the isolated MW snapshot which is extracted in a spherical shell with radii between 6 and 10~kpc from the Galactic center. The value of $\kappa_{\rm LMC}$ is specified in the top left corner of each panel. The panels below each speed distribution plot show the ratio of the MW+LMC and the MW-only distributions, for all snapshots except the isolated MW. 
    }
    \label{fig: fv panel plot halo 49}
\end{figure}



Figure~\ref{fig: fv panel plot halo 49} demonstrates that the LMC impacts the high speed tail of the local DM speed distribution, not only at its pericenter approach and at the present day, but also $\sim 175$~Myr after the present day. The value of $\kappa_{\rm LMC}$ is largest at pericenter and decreases as the LMC moves further from the host galaxy. Similarly, the ratio of the MW+LMC and the MW-only speed distributions in the high speed tail is largest at pericenter and decreases for the present day and future snapshots. For all  snapshots other than the isolated MW  (where $\kappa_{\rm LMC}=0$), the DM originating from the LMC has a speed distribution that peaks at the high speed tail of the native DM distribution of the MW, having a modest yet important impact on the total DM speed distribution. This is  similar to what we find in general for the 15 MW-LMC analogues at pericenter, as shown in figure~\ref{fig: fv lmc candidates}. 

Notice that when halo 13 is re-simulated with finer snapshots, the phase-space distribution of the DM particles  is not the same as in the original halo 13, as we are not comparing the two simulations at exactly the same time. As mentioned in section~\ref{sec: selection criteria}, the average time between snapshots is $\sim 150$~Myr in the original simulation, and it is difficult to precisely obtain the snapshot for the present day or LMC's pericenter approach. Furthermore, the Solar region for the best fit Sun's position is different in the original and the re-simulated halo 13, and this has a significant impact on the local DM velocity distribution. In particular, for the original halo 13, the cosine angles (eq.~\eqref{eq: cosine angles}) for the best fit Sun's position are $\cos \alpha=-0.796$ and $\cos\beta=-0.090$. Although these particular angles minimize the sum of the squared differences with their observed values, the value of $\cos\beta$ is very different than its observed value (as given in eq.~\eqref{eq: best fit cosine angles}), and it is therefore difficult to obtain a precise match to the observed Sun's position in the original halo 13. However, for the re-simulated halo 13, we obtain a much better match to the observed Sun's position (e.g.~$\cos \alpha=-0.995$ and $\cos\beta=-0.656$ for the best fit Sun's position at the present day snapshot). As a result, the speed distribution of DM particles from the LMC in the Solar region peaks at a noticeably higher speed in the re-simulated halo 13 compared to the original halo 13. Finally, there may also be a small variation in the phase-space distribution induced by the stochasticity of the baryonic physics model, which could lead to a slightly different evolution of the gravitational potential in the re-simulated halo. Hence, the local DM speed distributions and the values of  $\kappa_{\rm LMC}$  are also different between figures~\ref{fig: fv lmc candidates} and \ref{fig: fv panel plot halo 49} for halo 13. 

Our results in general confirm those presented in refs.~\cite{Donaldson:2021byu} and \cite{Besla:2019xbx}, which found that the small fraction of DM particles originating from the LMC in the Solar neighborhood (e.g.~$\sim 0.2\%$ in ref.~\cite{Besla:2019xbx}) dominates the high speed tail of the local DM speed distribution, in a suite of idealized simulations. Nevertheless, we note that the important effect of halo-to-halo variation in the results of our cosmological simulations cannot be overlooked.



